% !TEX root = ../main.tex
\section{Problema 5: Ubicación de Personas en una Reunión}

\subsection{Descripción del Problema y Abstracción CSP}
% Descripción: organización de N personas en una fila para una fotografía respetando sus preferencias.
% Preferencias:
% 1. next(A, B): la persona A debe estar junto a la persona B.
% 2. separate(A, B): la persona A debe estar separada de la persona B (no contiguas).
% 3. distance(A, B, M): las personas A y B deben estar separadas como máximo por M personas.

\subsubsection{Parámetros y Estructura de Entrada}
% Detallar los arreglos de entrada:
% personas, next (matriz k x 2), separate (matriz s x 2), distance (matriz w x 3).

\subsubsection{Variables de Decisión y Dominios}
% Explicar el modelado de variables (ej. pos[p]: posición de la persona p en la fila, o fila[i]: persona en la posición i).
% Justificar el uso de un modelo dual con canalización (inverse) si aplica.

\subsubsection{Restricciones del Modelo}
% Formalizar en lenguaje matemático:
% 1. Cada persona ocupa una única posición (alldifferent).
% 2. Restricción next(A, B): |pos[A] - pos[B]| = 1.
% 3. Restricción separate(A, B): |pos[A] - pos[B]| > 1.
% 4. Restricción distance(A, B, M): |pos[A] - pos[B]| - 1 <= M.

\subsection{Detalles de Implementación (\texttt{reunion.mzn})}

\subsubsection{Restricciones Redundantes y Rompimiento de Simetrías}
% Rompimiento de simetría de inversión de la fila (reflejo izquierda-derecha).
% Explicar qué simetría se rompe, cómo se implementa y cómo se recuperarían soluciones si se deseara.

\subsection{Estrategias de Búsqueda y Distribución}
% Anotaciones de búsqueda exploradas para la asignación de posiciones.

\subsection{Pruebas y Análisis Experimental}
% Presentar pruebas con diferentes tamaños de fila y conjuntos de restricciones (ej. instancias factibles e infactibles).

\begin{table}[htbp]
    \centering
    \begin{tabular}{lcccc}
        \toprule
        Instancia ($N$) & Estrategia & Tiempo (s) & Nodos explorados & Nodos fallidos \\
        \midrule
        Instancia base ($N=8$) & Default & -- & -- & -- \\
        Instancia base ($N=8$) & Custom & -- & -- & -- \\
        Instancia grande ($N=20$) & Default & -- & -- & -- \\
        \bottomrule
    \end{tabular}
    \caption{Resultados experimentales para el problema de la reunión.}
    \label{tab:reunion_resultados}
\end{table}

\subsection{Conclusiones del Ejercicio}
% Discusión de los resultados y efectividad del rompimiento de simetrías.
